\documentclass[11pt]{article}
\usepackage[margin=1in]{geometry}
\usepackage{amsmath,amssymb}
\usepackage{booktabs}
\newcommand{\relu}{\operatorname{ReLU}}
\newcommand{\sm}{\operatorname{softmax}}

\title{Implementing a \texttt{log} read-out with a ReLU feed-forward layer}
\date{}

\begin{document}
\maketitle
\vspace{-2.5em}

\section{Setup}
A block of the residual stream holds the in-context count vector
$\mathbf c=(c_1,\dots,c_V)$, the empirical next-token frequencies, with
$c_i\ge 0$ and $\sum_i c_i\le 1$. The feed-forward block and read-out are
\[
\mathrm{FF}(\mathbf s)=\mathbf s+W_2\,\relu(W_1\mathbf s+\mathbf b_1),
\qquad
\boldsymbol\ell = W_o\,\mathbf s+\mathbf b_o,
\qquad
\mathbf p=\sm(\boldsymbol\ell).
\]

\section{Goal}
The best count-based predictor is proportional to the (floored) counts,
\[
p_i=\frac{c_i+\beta_i}{\sum_j (c_j+\beta_j)} \qquad(\beta_i>0\ \text{a small floor}).
\]
Since $\sm(\log\mathbf u)_i=u_i/\sum_j u_j$, the soft-max produces exactly this \emph{iff}
\begin{equation}\label{eq:goal}
\boxed{\ \ell_i=\log(c_i+\beta_i)\ }.
\end{equation}
So the feed-forward layer must compute $\log(c_i+\beta_i)$ --- and we build it from ReLU only.

\section{ReLU makes piecewise-linear functions}
$\relu(c-x)$ is a hinge: $0$ for $c<x$, slope $1$ for $c>x$. Pick breakpoints
$0=x_0<x_1<\dots<x_{K-1}$ and let $m_j$ be the slope wanted on segment $[x_j,x_{j+1}]$.
Then
\begin{equation}\label{eq:pl}
\hat h(c)=\log\beta+\underbrace{m_0}_{\text{slope}}\relu(c-x_0)
+\sum_{j\ge 1}\underbrace{(m_j-m_{j-1})}_{\text{slope change}}\relu(c-x_j).
\end{equation}
On segment $k$ the slopes telescope to $m_0+\sum_{j=1}^{k}(m_j-m_{j-1})=m_k$, so each
piece has the chosen slope and the curve is continuous.

To approximate $\log(c+\beta)$, set $m_j$ to its secant slopes:
\begin{equation}\label{eq:slopes}
m_j=\frac{\log(x_{j+1}+\beta)-\log(x_j+\beta)}{x_{j+1}-x_j}.
\end{equation}
$\log$ is concave $\Rightarrow$ the $m_j$ decrease $\Rightarrow$ the first coefficient is
positive and the rest are negative. Error $\to 0$ as the number of knots $K$ grows.

\section{The weights}
One hidden unit per (coordinate $i$, knot $k$):
\[
W_1[u,\,\text{index of }c_i]=1,\qquad
b_1[u]=-x_k,\qquad
W_2[\text{out }i,\,u]=
\begin{cases}m_0 & k=0\\ m_k-m_{k-1} & k\ge 1\end{cases},
\qquad
b_o[i]=\log\beta_i .
\]
Then the unit outputs $\relu(c_i-x_k)$, the output coordinate accumulates
$\sum_k(\cdot)\,\relu(c_i-x_k)$, and adding $b_o[i]$ gives
$\ell_i=\hat h(c_i)\approx\log(c_i+\beta_i)$. Only $W_1$ ``selects'' a coordinate (a
single $1$ per row); $\log$ is used \emph{only} to precompute the fixed numbers
$m_j$, never on the data.

\section{Numbers ($\beta=0.01$, $K=7$, $6$ ReLU units per coordinate)}
\[
x_k:\ 0,\ 0.010,\ 0.025,\ 0.063,\ 0.159,\ 0.398,\ 1.0
\qquad(\text{knots, dense near }0).
\]
\begin{center}
\begin{tabular}{cccc}
\toprule
unit $u$ & ReLU bias $b_1=-x_k$ & $W_2$ coeff & meaning\\
\midrule
$0$ & $0.000$  & $+69.32$ & initial slope $m_0$\\
$1$ & $-0.010$ & $-32.08$ & bend down\\
$2$ & $-0.025$ & $-17.94$ & bend down\\
$3$ & $-0.063$ & $-10.55$ & bend down\\
$4$ & $-0.159$ & $-5.06$  & bend down\\
$5$ & $-0.398$ & $-2.19$  & bend down\\
\bottomrule
\end{tabular}
\end{center}
Read-out bias $b_o=\log(0.01)=-4.605$. With $V=20$ coordinates this uses
$20\times 6=120$ of the $128$ hidden units.

\section{Result}
\[
\sm(\boldsymbol\ell)_i
=\frac{e^{\log(c_i+\beta_i)}}{\sum_j e^{\log(c_j+\beta_j)}}
=\frac{c_i+\beta_i}{\sum_j(c_j+\beta_j)} .
\]
A pure-ReLU feed-forward layer, with hand-set weights, reproduces the divide-by-sum
(count) predictor: the $\log$ inverts the soft-max exponential. Measured cross-entropy
matches the exact-$\log$ read-out to $\approx 0.01$ nats at $K=7$, shrinking as $K$ grows.
\end{document}
